\scanpage{73}
% loc: 1–3
combinatoire $I_e$ sur $E^*/{\pm1}$ : structure de \deleted{graphe} \unsure{$k$|$\Delta$|bi}-polyèdre \uline{orienté}. Je rappelle que cette structure consiste en

% loc: 4–5
a) Structure de graphe : $\{\dot x,\dot y\}$ forment une arête \deleted{\unsure{si}} sss $\overbrace{x\wedge y}=e$ ou $-e$ ;

% loc: 6–8
$\longmapsto$ structure triangulée $\{\unsure{x|\dot x},y,z\}$ est un triangle ssi $\{\dot x,\dot y\}$, $\{\dot y,\dot z\}$, $\{\dot z,\dot x\}$ sont des arêtes.

% loc: 8–9, NB between a and b
\inserted{\uline{NB} a) Chaque arête est contenue dans exactement $\underline{2}$ triangles b) les triangles sont les triplets de la forme $\{\dot x,\dot y,\dot z\}$, avec $x+y+z=0$,}

% loc: 9–13
b) Chaque triangle est \uline{orienté} -- \deleted{(l'arête $\{\dot x,\dot y\}$ est directe (N} l'arête $\{\dot x,\dot y\}$ du triangle défini par $x,y,z$ avec $x\unsure{-|+}y+z=0$, \deleted{et $x\wedge y$} est \uline{directe}, sss $x\wedge y=e$, d'où $y\wedge z=e$, $z\wedge x=e$

% loc: 13–14, interlinear
\inserted{De telle façon que les deux or. d'une arête induites par celles des triangles adjacents soient opposées}

% loc: 15–24, lower part
c) L'``ens des arêtes et triangles incidents à une arête, qui est un contour combinatoire (\unsure{pas|non|un} néc. connexe) sur lequel $\mathbb Z$ opère (de par les données structurelles précédentes) est muni'' de plus d'une opération de $k$ dessus, prolongeant celle de $\mathbb Z$, et pour laquelle l'ens des arêtes \inserted{($\longmapsto$ sommets)} et l'ens des \deleted{triangles} ($\longmapsto$ arêtes) sont des \inserted{$\unsure{k|\lambda|\mathbb Z}$-torseurs}.
% archive sheet number: 72

